\documentclass[11pt]{article}
\usepackage{setspace}
\usepackage{geometry}
\geometry{margin=1in}
\setstretch{1.15}

\begin{document}

\noindent
\textbf{Application for Postdoctoral Researcher — Sustainable Consumer behavior Across Welfare Groups} \\[1mm]

\noindent Dear Prof.\ Granato and Members of the Selection Committee, \\[1mm]

I am writing to apply for the Postdoctoral Researcher position on \textit{Sustainable Consumer behavior Across Welfare Groups} at TU Delft. I am completing my Ph.D.\ in Economics at Bielefeld University and Università Cattolica del Sacro Cuore through the Marie Skłodowska–Curie EPOC programme. My research combines behavioral theory and computational modeling to study how heterogeneous individuals form beliefs, respond to incentives, and adapt behavior.

My background includes empirical studies in Experimental Economics. In my Master's coursework, I conducted an online survey with two colleagues to study the relationship between biased self-beliefs (over/underconfidence) and procrastination. During my Ph.D., I kept this interest by modeling human-like heuristics in computational simulations to analyze firm competition under bounded rationality and social influence. I am highly motivated to translate these modeling insights into experimental and survey-based research with human participants.
The core of my dissertation investigates behavioral responses to new technologies in environments characterized by informational frictions and socio-economic heterogeneity. Although focused on digital markets, the underlying mechanisms—preference heterogeneity, social influence, belief updating, and behavioral adaptation—map naturally onto sustainable consumption contexts. 
% This experience equips me to examine how psychological and socio-economic factors differ across welfare groups and how such heterogeneity shapes the effectiveness of behavioral interventions.

Beyond my dissertation, I have contributed to projects demonstrating my ability to work with modern quantitative methods and large-scale behavioral data. \textit{Polytope Fraud Theory} applies machine-learning techniques to detect behavioral anomalies in financial datasets, while a second project develops agentic workflows integrating large-language-model agents into behavioral analysis pipelines. These projects reflect strong analytical skills and practical experience with R, Python, and Mathematica. I am especially drawn to this position because of group’s expertise in sustainable behavior, de-consumption, and intervention design. Your work on affective and social determinants of consumer choices strongly resonates with my interest in modeling behavioral differences across socio-economic strata. I would welcome the opportunity to contribute to quantitative data collection, experimental design, survey analysis, and the development of differentiated behavioral strategies for promoting sustainable consumption.
At TU Delft, I hope to contribute to designing and analysing experiments and surveys on pro-environmental behavior across welfare groups. I plan to develop tailored intervention strategies based on behavioral heterogeneity and to integrate experimental findings with behavioral modeling for policy guidance.

Thank you so much for considering my application! I would welcome the opportunity to contribute my expertise to the faculty and I look forward to discussing how my research can support the development of sustainable consumption strategies.

% \vspace{6mm}

\noindent Sincerely,\\[10mm]
\noindent Zhongli Wang

\end{document}
